\documentclass[11pt,a4paper]{article}
\usepackage[margin=25mm]{geometry}
\usepackage{booktabs,tabularx,array}
\usepackage{amsmath,amssymb}
\usepackage{hyperref}
\usepackage{enumitem}
\usepackage{url}
\title{生成映像制作における構成要素の限界貢献と費用対効果\\
\large Software Factory・R2V・TPSに基づくアブレーション研究計画}
\author{Research proposal --- bonsai/r2v}
\date{\today}
\begin{document}
\maketitle

\begin{abstract}
本研究は、Reference-to-Video (R2V) 映像制作において、棒人間／ポーズ制御とJEV型の評価・修正ループが、品質と採用ショットあたり総費用に与える影響を検証する。初期実験では比較条件を3つに限定し、YAMLは全条件で実行条件と結果を記録する共通基盤として使用する。LoRA、追加制御アダプター、TPS型品質ゲートは初期実験から除外し、結果に応じて後続研究へ送る。目的は全要素の導入を正当化することではなく、寄与の小さい要素を除外し、最小コストで必要品質を満たす構成を特定することである。現段階では研究計画であり、改善効果や企業一般の課題を実証済みとは主張しない。
\end{abstract}

\section{研究背景}
生成映像は、単一クリップの視覚的品質に加えて、人物・衣装・物体の同一性、動作の明瞭さ、時間的一貫性、指示への適合を同時に満たす必要がある。既存の評価研究は、動画品質を複数の次元に分解し、自動指標と人間の判断の関係を調べている \cite{huang2024vbench,liu2024evalcrafter,han2025videobench,podstawski2025fcb}。長尺の物語生成では、キャラクターや場面の一貫性も重要な研究課題として整理されている \cite{elmoghany2025longvid}。

本研究の焦点は、新しい動画生成モデル自体を提案することではなく、既存の生成能力を制作工程に組み込み、品質と費用を管理する方法にある。Software Factoryの仕様・履歴管理は実験の共通基盤として採用するが、その独立した効果は初期実験では主張しない。TPS型品質ゲートも初期実験の独立変数から外す。

\section{企業の課題仮説}
制作現場で問題になり得る項目を、検証対象となる仮説として定義する。
\begin{itemize}[nosep]
\item キャラクターの外見や衣装がショット間で変化する。
\item 指定したポーズ・動作・画角が生成結果に反映されない。
\item 不合格ショットの再生成・手修正が多く、採用ショットあたりの費用が読みにくい。
\item 生成条件や採用理由が記録されず、成功例を再現できない。
\item レビュー工程が遅く、失敗の原因が仕様・参照・生成・評価のどこにあるか特定しづらい。
\end{itemize}
これらは企業全体に共通する事実として断言しない。公開事例・既存調査・実務者インタビュー・PoCで確認し、確認できなかった課題は研究動機から除外または限定する。

\section{研究目的と研究質問}
\begin{description}[style=nextline,leftmargin=0pt]
\item[RQ1] 棒人間／ポーズ制御は、同じR2Vベースラインに比べて指定動作への一致を改善し、追加費用に見合うか。
\item[RQ2] ポーズ制御にJEVの評価・修正ループを加えると、初回合格率・手戻り・人間レビュー時間は改善するか。
\item[RQ3] 人間の最低品質条件を満たした上で、どの条件が採用ショットあたり総費用とリードタイムの面で最も効率的か。
\end{description}

\section{提案システム}
構成要素を次のように役割分離する。
\begin{description}[style=nextline,leftmargin=0pt]
\item[Reference] 参照画像、キャラクター仕様、権利・出典情報、シーン意図を管理する。
\item[Generation] R2Vモデル、プロンプト、seed、LoRA、制御アダプター、実行パラメータを記録する。
\item[Specification] YAMLでショットの意図、必須条件、禁止条件、参照、評価基準を構造化する。
\item[Eval / JEV] 命題、証拠、ルーブリック、判定、信頼度、失敗分類を記録する。自動評価器の判定は人間評価と照合する。
\item[Production System / TPS] 品質ゲート、異常検出、原因工程への差し戻し、手戻りと待ち時間の計測を行う。
\item[Experiment tracking] Kaggle Notebook等でデータ版、モデル版、条件、ログ、費用、評価結果を再集計可能にする。
\end{description}

\section{サーベイの方針}
サーベイは、単なる関連研究の列挙ではなく、実験変数と評価指標を決めるために行う。主要な調査軸は以下の通り。
\begin{enumerate}[nosep]
\item \textbf{動画生成評価：} VBenchは複数の品質次元を分離した評価スイートを提案する \cite{huang2024vbench}。EvalCrafterは生成動画の評価ベンチマークを扱う \cite{liu2024evalcrafter}。Video-Benchは人間評価との整合を重視する \cite{han2025videobench}。
\item \textbf{被写体・顔の一貫性：} Face Consistency Benchmarkは、生成動画における顔の一貫性を測定する枠組みを扱う \cite{podstawski2025fcb}。Subject-to-Videoの評価では、被写体の同一性や自然さに焦点を当てるベンチマークも報告されている \cite{yuan2025opens2v}。
\item \textbf{長尺・物語生成：} 長尺映像の構成、一貫性、映像品質を整理したサーベイを参照し、短いPoCの結果を長編映画全体へ過剰一般化しない \cite{elmoghany2025longvid}。
\item \textbf{制御と計算効率：} ControlNeXtは動画・画像生成の制御を効率化する設計を検討している \cite{peng2024controlnext}。本研究では特定方式の優位性を前提とせず、制御品質と実行コストを測定する。
\item \textbf{製造・ソフトウェア工程：} TPSおよびSoftware Factoryの文献は、品質ゲート、流れ、手戻り、工程間の責任分界を設計するために調査する。製造業の知見が生成映像制作へそのまま移植できるとは仮定せず、適用可能性を実験対象とする。
\end{enumerate}
各論文について、課題、データセット、比較条件、指標、人間評価の有無、計算コスト、再現可能性、ライセンスをサーベイ表に記録する。引用文献の主張は原論文本文で確認し、検索結果の要約だけを根拠に新規性を主張しない。

\section{実験設計：貢献しない要素を見抜く}
初期実験では変数を増やさない。YAML、同一モデル、評価ルーブリック、ログ形式は全条件で固定・共有し、実験変数は「ポーズ制御の有無」と「評価・修正ループの有無」に限定する。これにより、YAMLの効果やLoRAの効果まで同時に主張する混同を避ける。

\subsection{基本比較}
\begin{center}
\small
\begin{tabularx}{\linewidth}{lX}
\toprule
条件 & 構成 \\
\midrule
A & R2Vベースライン（YAMLログ基盤を共通使用） \\
B & A + 棒人間／ポーズ制御 \\
C & B + JEV型評価・修正ループ \\
\bottomrule
\end{tabularx}
\end{center}
初期実験ではA/B/Cのみを実施する。Cの改善が確認され、計算資源とデータが確保できた場合に限り、LoRAまたは追加アダプターを次段階の独立実験として検討する。TPSゲートは実際のレビュー工程を測定できる段階まで保留する。

\subsection{アブレーションと除外試験}
各要素 $i$ の限界寄与を、基準構成から当該要素だけを除いた条件との差で推定する。
\[
\Delta Q_i = Q(\text{full}) - Q(\text{full}\setminus i)
\]
\[
\Delta C_i = C(\text{full}) - C(\text{full}\setminus i)
\]
ここで $Q$ は事前に定義した品質指標、$C$ は採用ショットあたりの総費用（推論、学習・調整、評価、再生成、レビュー工数を含む）である。要素の導入による費用対効果は、品質改善を追加費用で割った値として補助的に報告する。ただし、品質指標が異なる尺度を持つ場合は、単純な比率だけで判断しない。

\[
\mathrm{CostPerAcceptedShot} =
\frac{\mathrm{GenerationCost}+\mathrm{TrainingCost}+\mathrm{EvaluationCost}+\mathrm{HumanReviewCost}}
{\mathrm{NumberOfAcceptedShots}}
\]
\[
\mathrm{FirstPassYield} =
\frac{\mathrm{FirstPassAcceptedShots}}{\mathrm{AllAttemptedShots}}
\]
要素を一つずつ追加する比較と、最終構成から一つずつ外す比較を併用する。前者は導入効果、後者は冗長性を検討する。複数要素の相互作用がある場合、単独比較だけでは結論を出さない。

\subsection{除外判断ルール}
\begin{itemize}[nosep]
\item \textbf{Keep：} 主要品質指標を実務上意味のある幅で改善し、その効果が反復試行でも確認され、追加コストを正当化できる。
\item \textbf{Drop：} 品質改善が事前に定めた最小実用効果を下回り、費用・時間・複雑さだけが増える。
\item \textbf{Conditional：} 特定カテゴリ（例：複雑な動作、複数人物）だけで効果がある。対象タスクに限定して使う。
\item \textbf{Inconclusive：} 試行数不足、評価者間の不一致、信頼区間の広さなどで結論できない。貢献なしと断定せず、追加データまたは実験の縮小を検討する。
\end{itemize}
統計的に有意でないことだけを理由に無効と断定しない。効果量、信頼区間、実用上の閾値、コストを併せて判断する。

\section{課題ショットとデータ選定}
課題は、(i) キャラクター同一性、(ii) ポーズ・動作、(iii) 構図・カメラ、(iv) 画風、(v) 時間的一貫性、(vi) 複数参照の結び付け、(vii) 遮蔽や交差動作などの難例に層別する。成功例だけを選ばず、失敗しやすいケースも含める。参照素材の出典、ライセンス、取得日、利用可能範囲を記録し、学習・調整用と最終評価用を分離する。

\section{Kaggle実験運用とコスト管理}
Kaggleは必須の実行基盤ではなく、Notebook、データ、実験ログを管理する候補環境と位置づける。まず少数の課題で環境確認（GPU・メモリ・推論時間・依存関係・ライセンス）を行い、ベースラインを成立させる。その後、情報量の高い比較を優先して条件を追加する。無条件の総当たりは避ける。
各実行について、モデル版、データ版、YAML版、プロンプト、seed（使用可能な場合）、パラメータ、実行時間、GPU時間、失敗理由、出力URI、評価結果を保存する。計算量が大きい学習は、推論時の制御実験で寄与が見込めた場合に限って着手する。

\section{評価指標}
\begin{itemize}[nosep]
\item 動作／ポーズ一致、キャラクター一貫性、時間的一貫性、指示適合、画質。
\item JEVと人間評価の一致度、評価者間一致度。
\item 初回合格率、手戻り率、再生成回数、レビュー時間、採用ショットあたり総費用。
\item 再現性：同一仕様・同一条件から得られる結果の変動。
\end{itemize}
主要指標と実用上の最小改善幅は実験前に固定する。自動評価器は誤判定や偏りを持ち得るため、人間の盲検評価を一部に用いる。結果を見てから課題を都合よく除外しない。

\section{実施計画}
\begin{enumerate}[nosep]
\item 文献サーベイと企業課題仮説の確認。
\item 課題ショット、ライセンス、ルーブリック、主要指標の事前固定。
\item ベースラインと実験ログの実装。
\item A/B/Cの3条件でポーズ制御とJEVループの寄与を比較する。
\item 人間の盲検評価とJEVを照合し、誤PASS率・レビュー時間・総費用を集計する。
\item 初期結果が正当化した場合に限り、LoRA、追加アダプター、TPSゲートを別の後続実験として評価する。
\item 結果・費用・失敗例の集計と再現確認。
\item 論文執筆。未測定値を埋めず、限界と否定的結果も報告する。
\end{enumerate}

\section{期待される貢献と限界}
期待する貢献は、(1) 生成映像制作の工程・仕様・評価を結ぶ再現可能な実験設計、(2) 品質改善と総費用を同時に測る構成要素別の寄与分析、(3) 寄与が乏しい構成要素を除外するための判断基準、(4) JEV評価と人間判断の一致度の検証である。実験前に、品質向上、コスト削減、企業への一般化、新規性を確定的に主張しない。

\section{参考文献}
\bibliographystyle{plain}
\bibliography{references}
\end{document}
